% 依赖项（必须放在主文档导言区，即 \begin{document} 之前）：
% \usepackage{tabularray}
% \UseTblrLibrary{varwidth}
% \usepackage{enumitem}
\begingroup
% \nopagebreak[4]
\enlargethispage{1cm}
\footnotesize
\centering

\DefTblrTemplate{conthead-text}{default}{（续表）}
\DefTblrTemplate{contfoot-text}{default}{续下页}

\begin{longtblr}[
  caption = {开源基础底座、云端企业智能体与自动化攻防评测体系},
  label   = {tab:industry_academia_agent_security},
]{
  width   = \linewidth,
  colspec = {
    |Q[wd=0.07\linewidth,l,m]
    |Q[wd=0.12\linewidth,l,m]
    |X[1.00,l,m]|
  },
  cells   = {l,m},
  colsep  = 3pt,
  row{1}  = {font=\bfseries},
  rowhead = 1,
  rowsep  = 5pt,
  measure = vbox,
  stretch = -1,
}
\hline
研究方向 & 企业/机构 + 核心学者 & 自研项目、产学合作、商用安全产品 \\
\hline

\SetCell[r=2]{l,m} 开源基础大模型原生智能体安全底座
& Meta FAIR；Yann LeCun, Joelle Pineau
& \begin{itemize}[leftmargin=*,topsep=0pt,partopsep=0pt,itemsep=1pt,parsep=0pt]
\item \textbf{世界模型安全底座：}以JEPA、Image-JEPA、Video-JEPA、V-JEPA 2构成学术路线演进，支撑多模态智能体的状态表征、感知预测和行动后果建模~\cite{nyuLeCun2026,vjepa2025}；
\item \textbf{开源安全护栏体系：}公开Llama Guard、Prompt Guard和LlamaFirewall，其中LlamaFirewall定位为面向安全AI Agent的开源护栏系统~\cite{metaLlamaFirewall2025}；
\item \textbf{团队组织演进：}Joelle Pineau卸任FAIR负责人、LeCun离开Meta后，相关研发由公司新组建的Meta Superintelligence Labs承接~\cite{pineauDeparture2025,lecunDeparture2025}。
\end{itemize} \\
\cline{2-3}
& 纽约大学、AMI Labs；Yann LeCun
& \begin{itemize}[leftmargin=*,topsep=0pt,partopsep=0pt,itemsep=1pt,parsep=0pt]
\item \textbf{JEPA架构产业化：}LeCun离职后创办AMI Labs推动JEPA架构商业化，首批目标行业为医疗、机器人和工业自动化~\cite{amiLabsFunding2026,lecunMITTechReview2026}；
\item \textbf{理论与基准验证：}2026年发表的JEPA相关论文证明该架构在特定条件下能够恢复真实世界底层结构，同时用配套基准指出当前模型面对细微视觉变化时仍较脆弱~\cite{lejepaIdentifiability2026}。
\end{itemize} \\
\hline

\SetCell[r=3]{l,m} 云端原生企业级大模型智能体研发与风控
& Google DeepMind；Demis Hassabis
& \begin{itemize}[leftmargin=*,topsep=0pt,partopsep=0pt,itemsep=1pt,parsep=0pt]
\item \textbf{多智能体协同基础理论：}AlphaStar、RoboCat等研究成果支撑多智能体协同博弈和机器人任务规划的基础理论；
\item \textbf{企业级智能体产品落地：}相关研究能力被工程化为Gemini Enterprise Agent Platform和Managed Agents，落实到API注册、IAM权限管控和运行日志审计~\cite{googleGeminiEnterpriseAgentPlatform2026,googleADK2025}；
\item \textbf{科研智能体协作平台：}AI Co-Scientist基于Gemini构建，由一组专业化Agent协同完成科学假设的生成、辩论、排序和迭代~\cite{deepmindCoScientist2025,gottweisAICoScientist2025}。
\end{itemize} \\
\cline{2-3}
& Google Research；Jeff Dean
& \begin{itemize}[leftmargin=*,topsep=0pt,partopsep=0pt,itemsep=1pt,parsep=0pt]
\item \textbf{云端Agent风控机制：}分布式大规模机器学习与系统工程能力支撑Gemini Enterprise Agent Platform的API注册、IAM权限管控和运行日志审计~\cite{googleADK2025,googleGeminiEnterpriseAgentPlatform2026}；
\item \textbf{托管式智能体治理：}Managed Agents将全链路Tracing审计、IAM权限管控等系统工程理论嵌入托管式智能体服务，对应企业Agent运行时风控需求~\cite{googleADK2025}。
\end{itemize} \\
\cline{2-3}
& 斯坦福大学医学院、MIT；Gary Peltz, Ritu Raman
& \begin{itemize}[leftmargin=*,topsep=0pt,partopsep=0pt,itemsep=1pt,parsep=0pt]
\item \textbf{产学研安全验证：}AI Co-Scientist于2025年2月首发时联合全球100多个科研机构共同开发和测试，Gary Peltz、Ritu Raman公开验证其科研任务能力~\cite{deepmindCoScientist2025}；
\item \textbf{科研落地与企业试用：}Co-Scientist已产出抗菌素耐药性、植物免疫等湿实验验证成果，并被第一三共、拜耳作物科学等企业内部试用~\cite{labcriticsCoScientist2026}。
\end{itemize} \\
\hline

\SetCell[r=2]{l,m} 智能体攻防评测基准与自动化红队体系
& 苏黎世联邦理工SPY Lab；Florian Tramèr
& \begin{itemize}[leftmargin=*,topsep=0pt,partopsep=0pt,itemsep=1pt,parsep=0pt]
\item \textbf{自动化攻防基准：}AgentDojo将间接提示注入、工具劫持等标准化攻防任务集统一到动态评测环境中，已由ETH SPY Lab开源部署、~\cite{agentdojoGithub2026}；
\item \textbf{基准复现实验：}AgentDojo使邮件、日历、文件和工具调用场景中的提示注入攻击具备可复现、可比较的实验对象。
\end{itemize} \\
\cline{2-3}
& CMU, Gray Swan AI；Zico Kolter, Matt Fredrikson, Andy Zou, Dawn Song
& \begin{itemize}[leftmargin=*,topsep=0pt,partopsep=0pt,itemsep=1pt,parsep=0pt]
\item \textbf{工业级红队平台：}Gray Swan AI将模型鲁棒性、神经网络形式化验证理论转化为检测能力，代表性产品包括红队测试平台Cygnal和对抗性红队工具Shade~\cite{grayswanSeriesA2026}；
\item \textbf{提示注入鲁棒性评测：}Shade已被Anthropic用于评估模型在代码环境中对提示注入攻击的鲁棒性，Gray Swan的间接提示注入研究也被Anthropic Claude Mythos的system card引用为权威研究依据~\cite{latentspaceGrayswan2026}；
\item \textbf{模型风险治理嵌入：}Kolter已出任OpenAI董事会成员并兼任安全与安保委员会主席，使自动化红队研究与模型风险治理进入更高层级的产业治理流程~\cite{kolterWikipedia2026}。
\end{itemize} \\
\hline
\end{longtblr}
\endgroup
